% SEGMENT OLP-0461-S01
% Part: normal-modal-logic
% Chapter: tableaux
% Section: introduction

\documentclass[../../../include/open-logic-section]{subfiles}

\begin{document}

\olfileid{nml}{tab}{int}

\olsection{அறிமுகம்}

!!^{tableau}s என்பவை !!{signed formula}s இன் கீழ்நோக்கிக் கிளைக்கும் சில
மரவுருக்கள்; அதாவது, ஒரு மெய்மதிப்புக் குறியும் ($\True$ அல்லது $\False$)
!!a{sentence} உம் சேர்ந்த சோடிகளின் மரவுருக்கள்:
\[
\sFmla{\True}{!A} \text{ அல்லது } \sFmla{\False}{!A}.
\]
!!^a{tableau} சில \emph{எடுகோள்களுடன்} தொடங்குகிறது. அதற்குப் பிந்தைய
ஒவ்வோர் !!{signed formula} உம் அனுமான விதிகளில் ஒன்றைப் பயன்படுத்துவதால்
உருவாக்கப்படுகிறது. சில அனுமான விதிகள் மரவுருவின் ஒரு முனையில் ஒன்று அல்லது
அதற்கு மேற்பட்ட !!{signed formula}s ஐச் சேர்க்கின்றன; மற்றவை இரண்டு புதிய
முனைகளைச் சேர்த்து இரண்டு கிளைகளை உருவாக்குகின்றன. ஒரு விதி பயன்படுத்தப்படும்
!!{signed formula}வில் உள்ள !!{formula}வைவிடச் சிக்கல் குறைந்த
வாய்பாடுகளைக் கொண்ட !!{signed formula}s ஐ விதிகள் விளைவிக்கின்றன. ஒரு
கிளையில் $\sFmla{\True}{!A}$, $\sFmla{\False}{!A}$ ஆகிய இரண்டும் இருந்தால்,
அந்தக் கிளை \emph{மூடியது} என்கிறோம். !!a{tableau}வின் ஒவ்வொரு கிளையும்
மூடியிருந்தால், அந்த முழு !!{tableau}வும் மூடியது. மூடிய !!a{tableau},
அந்த !!{tableau}வைத் தொடங்கப் பயன்படுத்திய !!{signed formula}s இன் கணம்
நிறைவுறத்தக்கதல்ல
என்பதைக் காட்டும் !!a{derivation} ஆகும். இதைக் கொண்டு $\Proves$ உறவை
வரையறுக்கலாம்: $\Gamma \Proves !A$ அப்போதும் அப்போதுமட்டுமே,
$\Gamma_0 = \{!B_1, \dots, !B_n\} \subseteq \Gamma$ என்ற ஏதோவொரு
முடிவுறு கணம் இருந்து, பின்வரும் எடுகோள்களுக்கு ஒரு மூடிய !!{tableau}
இருக்கும்:
\[
\{\sFmla{\False}{!A}, \sFmla{\True}{!B_1}, \dots, \sFmla{\True}{!B_n}\}.
\]

வாய்ப்புநிலைத் தருக்கங்களுக்கு, !!{signed formula} என்ற கருத்தை விரிவுபடுத்தி,
~\iftag{prvBox}{$\Box$\iftag{prvDiamond}{ மற்றும்
    $\Diamond$}}{$\Diamond$} ஆகியவற்றைக் கையாளும் விதிகளையும் சேர்க்க
வேண்டும். ஒரு குறியுடன் ($\True$ அல்லது $\False$) கூடவே, வாய்ப்புநிலை
!!{tableau}s இல் உள்ள !!{formula}s க்கு \emph{முன்னொட்டுகள்}~$\sigma$ உம்
உண்டு. முன்னொட்டுகள் நேர்ம முழுக்களின் வெறுமையற்ற தொடர்கள்; அதாவது,
$\sigma \in (\PosInt)^* \setminus \{\emptyseq\}$. அத்தகைய முன்னொட்டுகளைச்
சுற்றியுள்ள $\tuple{\ }$ இல்லாமல் எழுதுகிறோம்; தனித்தனி !!{element}s ஐப்
பிரிக்க~$,$ களுக்குப் பதில்~$.$ களைப் பயன்படுத்துகிறோம். $\sigma$
ஒரு முன்னொட்டு எனில், $\sigma.n$ என்பது $\sigma \concat \tuple{n}$;
எடுத்துக்காட்டாக, $\sigma = 1.2.1$ எனில், $\sigma.3$ என்பது $1.2.1.3$.
ஆகவே, எடுத்துக்காட்டாக,
\[
\sFmla{\True}{\Box !A \lif !A}[1.2]
\]
என்பது ஒரு \emph{முன்னொட்டுள்ள !!{signed formula}} (அல்லது சுருக்கமாக ஒரு
\emph{முன்னொட்டுள்ள !!{formula}}).

உள்ளுணர்வின்படி, !!a{tableau}வின் ஒரு கிளையிலுள்ள !!{formula}s ஐ
நிறைவுறுத்தக்கூடிய மாதிரியின் ஓர் உலகிற்கு முன்னொட்டு பெயரிடுகிறது;
$\sigma$ ஓர் உலகிற்குப் பெயரிடுகிறது எனில், $\sigma.n$ ஆனது~$\sigma$
பெயரிடும் அந்த உலகிலிருந்து அணுகத்தக்க ஓர் உலகிற்குப் பெயரிடுகிறது.

\end{document}
